% =====================================================================
%  SUPPLEMENTARY MATERIAL — RMAT-2026-0601, round-1 revision
%  Accompanies paper_PGAE_oneclass.tex. Isolated working copy under
%  review vong 1/response/manuscript/; does not touch the original
%  manuscript/ tree.
% =====================================================================
\documentclass[11pt,a4paper]{article}

\usepackage[utf8]{inputenc}
\usepackage[T1]{fontenc}
\usepackage[margin=2.4cm]{geometry}
\usepackage{amsmath,amssymb}
\usepackage{booktabs}
\usepackage{longtable}
\usepackage[hidelinks]{hyperref}
\usepackage{xcolor}
\usepackage{caption}
\usepackage{authblk}

\renewcommand{\thesection}{S\arabic{section}}
\renewcommand{\thetable}{S\arabic{table}}
\renewcommand{\thefigure}{S\arabic{figure}}
\renewcommand{\theequation}{S\arabic{equation}}

\title{\textbf{Supplementary Material}\\[4pt]
\large Physics-Guided Autoencoders with Unsupervised Feature
Attention for One-Class Structural Damage Detection}

\author[1,2]{Chung Truong Thanh}
\author[1]{Hoa Tran Ngoc\thanks{Corresponding author: \texttt{ngochoa@utc.edu.vn}}}
\author[1]{Thanh Bui Tien}
\affil[1]{University of Transport and Communications, No.\ 3 Cau Giay Street, Hanoi, Vietnam}
\affil[2]{Nha Trang University, 02 Nguyen Dinh Chieu, Khanh Hoa, Vietnam}
\date{}

\begin{document}
\maketitle
\vspace{-2em}

\section{Sensitivity to constant features}
\label{sec:constant}

Section 3.4 of the main text notes that a small number of features are
constant across every window and receive the largest attention weight for
that reason alone, since a feature with zero variance is reconstructed with
zero error. Because a constant feature cannot distinguish healthy from
damaged windows, that section argues it adds nothing to the detection score.
Here we test that argument directly rather than leave it as a matter of
formula.

We identify every feature with essentially zero variance across the
$256$-window healthy training pool, using the same threshold as the
constant-feature check reported in the code release. Seventy of the $906$
features meet this criterion. We remove them entirely, retrain PGAE-A on the
remaining $836$ features, and repeat the automatic $\beta$ selection and the
ten-seed evaluation protocol used for Table 5 of the main text, so the two
conditions can be compared directly.

Table~\ref{tab:constant} reports the result. Both $\beta$ and the
participation ratio fall once the constant features are removed, $\beta$
from $0.065$ to $0.048$ and the participation ratio from $282$ to $218$
features. This shift is expected: the constant features previously occupied
several of the lowest-residual positions in the ranking that $\beta$ is fit
against, and removing them changes the shape of that ranking. AUC, in
contrast, is essentially unchanged, from $0.795$ to $0.789$, and the AUC
values from the two conditions overlap almost entirely across the same ten
seeds.

The constant features therefore behave exactly as Section 3.4 of the main
text predicts: they shift where the attention mechanism places its internal
cutoff, but they do not affect what the mechanism detects. Retaining them in
the feature bank has no measurable effect on detection performance.

\begin{table}[htbp]
\centering
\caption{Effect of removing the $70$ constant features on $\beta$, the
participation ratio, and AUC (ten seeds; same evaluation set as Table 5 of
the main text).}
\label{tab:constant}
\small
\begin{tabular}{@{}lcccc@{}}
\toprule
\textbf{Condition} & \textbf{Features} & \boldmath{$\beta$} & \textbf{Participation ratio} & \textbf{AUC} \\
\midrule
Full                          & $906$ & $0.065 \pm 0.007$ & $282 \pm 5$ & $0.795 \pm 0.014$ \\
Constant features removed     & $836$ & $0.048 \pm 0.005$ & $218 \pm 7$ & $0.789 \pm 0.016$ \\
\bottomrule
\end{tabular}
\end{table}

\section{Sensitivity to the calibration subset}
\label{sec:calibration}

Within the standard ten-seed protocol, the residual statistics that set
$\beta$ and the attention weights are computed from $X_{\text{train}}$, the
same rows used to fit the autoencoder by backpropagation. Because the input
is $906$-dimensional and the training split contains only around two hundred
windows, a small residual on $X_{\text{train}}$ could partly reflect the
network fitting those particular rows rather than a property of the healthy
state that generalizes. We test this directly by computing weights from
$X_{\text{val}}$ instead: the $20\%$ of the healthy training pool held out
from every gradient update, used only to decide when to stop training.
$X_{\text{val}}$ is disjoint from the evaluation set, so the comparison below
is not circular.

Table~\ref{tab:calibration} reports $\beta$ and AUC under both calibration
sources, using the same trained autoencoder and the same evaluation set per
seed. $\beta$ roughly doubles when computed from $X_{\text{val}}$, from
$0.065$ to $0.123$, consistent with the smaller sample available for
calibration, about fifty rows instead of two hundred. AUC, in contrast, is
unchanged: $0.795$ with self-calibration versus $0.796$ with independent
calibration, and a paired comparison across the same ten seeds shows no
significant difference. The two resulting weight vectors are themselves
highly correlated, $0.93$ on average.

Estimating the attention weights from the autoencoder's own training rows
therefore does not inflate detection performance relative to an independent
calibration subset. The shift in $\beta$ reflects the smaller sample size
available for calibration, not overfitting.

\begin{table}[htbp]
\centering
\caption{Detection AUC under self-calibration versus an independent
calibration subset (ten seeds; same evaluation set as Table 5 of the main
text).}
\label{tab:calibration}
\small
\begin{tabular}{@{}lccc@{}}
\toprule
\textbf{Calibration source} & \textbf{Rows} & \boldmath{$\beta$} & \textbf{AUC} \\
\midrule
$X_{\text{train}}$ (self-calibration, standard protocol) & $\approx 205$ & $0.065 \pm 0.007$ & $0.795 \pm 0.014$ \\
$X_{\text{val}}$ (independent calibration)                & $\approx 51$  & $0.123 \pm 0.015$ & $0.796 \pm 0.018$ \\
\bottomrule
\end{tabular}
\end{table}

\section{Sensitivity to feature redundancy}
\label{sec:redundancy}

The feature bank is deliberately redundant: the same physical event can leave
a trace in several descriptors from the same channel or in the same
descriptor across several channels. A feature could therefore receive a
large attention weight simply because it duplicates information already
carried by other features, rather than because it captures something
specific to the healthy structure. Here we test how much of the
attention-weight pattern this redundancy explains.

Section~\ref{sec:constant} sets aside the $70$ constant features, which are
reconstructed perfectly for a trivial reason; the remaining $836$
non-constant features are the subject here. For each of them, we compute two
independent quantities. The first is the feature's average attention weight across the
ten seeds of the standard protocol. The second is a redundancy score: the
largest absolute Pearson correlation between that feature and any other of
the $836$, computed on the $32$-window healthy test pool, which is held out
from training and from attention calibration in every seed. Because this
redundancy score never touches the data used to fit the autoencoder or the
attention weights, a relationship between the two quantities cannot be an
artifact of the same data being reused twice.

The two quantities are strongly related. Across all $836$ features, the
Spearman rank correlation between attention weight and redundancy is
$\rho=0.611$ ($p=8.8\times10^{-87}$): features that duplicate another
feature in the bank systematically receive larger weight. The relationship
is sharpest at the top of the weight distribution. Table~\ref{tab:redundancy}
splits the $836$ features into the twenty receiving the largest average
weight and the remaining $816$. The top twenty carry, on average, $5.5$
times the mean weight of a non-constant feature, and sixteen of them
($80\%$) correlate above $0.99$ with at least one other feature, against
$40\%$ among the remaining features.

Feature redundancy therefore accounts for a substantial part of the
attention-weight pattern, as expected from the redundancy of the feature
bank itself. This is consistent with reading the attention weights as an
operational reconstruction-reliability ranking rather than as evidence that
the structure specifically controls whichever physical quantity a
high-weight feature happens to describe.

\begin{table}[htbp]
\centering
\caption{Redundancy of the $836$ non-constant features, split by attention
weight (ten-seed average weight; redundancy is the largest absolute
correlation with another feature, measured on the $32$-window healthy test
pool).}
\label{tab:redundancy}
\small
\begin{tabular}{@{}lccccc@{}}
\toprule
\textbf{Group} & \textbf{Features} & \textbf{Mean weight} & \textbf{Mean redund.} & \textbf{Median redund.} & \textbf{Redund.\ $>0.99$} \\
\midrule
Top $20$ by weight   & $20$  & $5.53\times$ mean & $0.932$ & $0.9999$ & $16/20$ ($80\%$) \\
Remaining            & $816$ & $0.89\times$ mean & $0.874$ & $0.965$  & $322/816$ ($39\%$) \\
All non-constant     & $836$ & $1.00\times$ mean & $0.876$ & $0.967$  & $338/836$ ($40\%$) \\
\bottomrule
\end{tabular}
\end{table}

\section{Leave-one-state-out validation}
\label{sec:loso}

The healthy training pool contains $256$ windows drawn from only eight
structural states, $32$ windows per state. Because windows from the same
state come from the same measurement campaign, they are more similar to each
other than windows from different states, so the random $80/20$
training/validation split used in the main text can place correlated
windows from the same state on both sides of the split. We replace it with
leave-one-state-out validation: each of the eight healthy states is held out
in turn as the validation set, with the remaining seven states, $224$
windows, as the training set. Attention weights are calibrated from the
training fold, matching the main text's standard protocol. The evaluation
set, $324$ damaged and $32$ healthy test-pool windows, is unchanged and
identical to Table 5 of the main text, so AUC is directly comparable.

Table~\ref{tab:loso} reports the result for each held-out state. The mean
AUC across the eight folds is $0.780 \pm 0.012$, compared with
$0.795 \pm 0.014$ under the random split. A Welch $t$-test comparing the
eight leave-one-state-out folds against the ten random-split seeds confirms
this $0.015$ difference is statistically significant, $p=0.031$. Every fold,
including the weakest, S4 at $0.762$, remains far above AE-plain's $0.643$.

The random split therefore gives a mildly optimistic estimate of detection
performance, consistent with correlated windows inflating the validation
signal used for early stopping. The magnitude of the effect is small
relative to the gap the paper's central claim rests on: attention improves
AUC over AE-plain by roughly $15$ points, and leave-one-state-out validation
reduces that improvement by about $1.5$ points, not by an amount that
changes the conclusion.

\begin{table}[htbp]
\centering
\caption{Detection AUC under leave-one-state-out validation, one fold per
healthy state (same evaluation set as Table 5 of the main text).}
\label{tab:loso}
\small
\begin{tabular}{@{}lccc@{}}
\toprule
\textbf{Held-out state} & \textbf{Training windows} & \boldmath{$\beta$} & \textbf{AUC} \\
\midrule
S1 & $224$ & $0.061$ & $0.775$ \\
S2 & $224$ & $0.079$ & $0.782$ \\
S3 & $224$ & $0.057$ & $0.771$ \\
S4 & $224$ & $0.088$ & $0.762$ \\
S5 & $224$ & $0.071$ & $0.775$ \\
S6 & $224$ & $0.088$ & $0.782$ \\
S7 & $224$ & $0.044$ & $0.793$ \\
S8 & $224$ & $0.045$ & $0.801$ \\
\midrule
\textbf{Mean $\pm$ std} & & $0.067 \pm 0.017$ & $0.780 \pm 0.012$ \\
\bottomrule
\end{tabular}
\end{table}

\section{Cluster bootstrap confidence interval for AUC}
\label{sec:bootstrap}

The Hanley--McNeil standard error reported in the main text treats the test
windows as independent observations. In fact the $324$ damaged windows come
from only nine structural states, $36$ windows each, and the $32$ healthy
test-pool windows from eight states, four windows each; windows from the
same state share the same measurement campaign and physical condition and
are therefore correlated, not independent. We estimate the sampling
distribution of AUC directly with a bootstrap, using the scores from the
main text's representative run.

Two bootstrap variants are compared. The first resamples individual windows
with replacement, ignoring the state structure; it serves as a check against
the Hanley--McNeil formula under the same independence assumption the
formula makes. The second, a cluster or block bootstrap, resamples states
with replacement instead of windows: each replicate draws nine damaged
states and eight healthy states, with repetition allowed, and pools all
windows belonging to the drawn states before computing AUC. This is the
standard way to obtain a valid confidence interval when observations are
grouped into correlated clusters. Both variants use $5000$ replicates.

Table~\ref{tab:bootstrap} reports the result. Both bootstrap standard errors
exceed the Hanley--McNeil value, $0.055$ and $0.049$ against $0.032$: the
formula in the main text understates the true uncertainty in the reported
AUC. The block bootstrap is not wider than the window-level bootstrap in
this dataset; if anything it is slightly narrower. This is a property of
the Z24 partition rather than a sign that state dependence does not matter:
every damaged state contributes exactly $36$ windows and every healthy-test
state exactly $4$, and window-to-window variability within a single state is
already substantial, so resampling whole states does not add as much extra
spread as it would under an unbalanced design. The comparison that matters
is not block versus window-level bootstrap, but either against the
Hanley--McNeil formula: both nonparametric estimates agree that the true
$95\%$ interval for AUC is closer to $\pm0.10$ than to the $\pm0.06$
implied by $\pm 1.96$ times the Hanley--McNeil standard error.

\begin{table}[htbp]
\centering
\caption{Standard error and $95\%$ interval for AUC under three uncertainty
estimates, computed from the main text's representative run (point
AUC $=0.802$).}
\label{tab:bootstrap}
\small
\begin{tabular}{@{}lcc@{}}
\toprule
\textbf{Method} & \textbf{SE} & \textbf{95\% interval} \\
\midrule
Hanley--McNeil (main text Eq.~8, assumes independent windows) & $0.032$ & --- \\
Bootstrap, resampling individual windows ($5000$ reps)         & $0.055$ & $[0.685,\ 0.899]$ \\
Cluster bootstrap, resampling states ($5000$ reps)              & $0.049$ & $[0.705,\ 0.895]$ \\
\bottomrule
\end{tabular}
\end{table}

\section{Feature bank definitions and implementation settings}
\label{sec:featuredefs}

The $906$-feature bank comprises three groups applied to each of the five
channels (R1V, R2L, R2T, R2V, R3V) or, for transmissibility, each of the
four non-reference channels paired with R1V: $49$ time-domain descriptors
per channel ($245$ features), $74$ PSD-shape descriptors per channel
($370$ features), and the same $74$ descriptors applied to the four
transmissibility functions ($296$ features). This section lists every
descriptor. The exact implementation is the single source of truth and is
available in full at the repository named in the Data and Code Availability
statement of the main text, specifically \texttt{code/features/utils4.py}
(time-domain) and \texttt{code/features/run\_eda\_pipeline.py}, function
\texttt{frequency\_domain\_features} (PSD/transmissibility).

\subsection{Time-domain descriptors (49 per channel)}

Table~\ref{tab:tddefs} lists all $49$. Most are computed with the
\texttt{tsfresh} feature-calculator library \cite{christ2018}; entries
below name the corresponding tsfresh function where applicable. The autocorrelation-family
descriptors (autocorrelation, partial autocorrelation, c3, time-reversal
asymmetry) use a per-window lag of $\min(\lfloor n/100\rfloor, 3000)$
samples, $n$ the window length. Five descriptors---absolute energy,
absolute sum of changes, PSD entropy, FFT entropy, and cid-ce complexity---
are computed by disabled code paths in the released implementation and are
therefore identically zero for every window; one further pair, ``fraction of
zero values'' and ``count of zeros,'' duplicates the same tsfresh
calculator. Together, these account for part of the constant and redundant
columns discussed in Sections~\ref{sec:constant} and~\ref{sec:redundancy}.

\begin{longtable}{@{}p{0.5cm}p{5.8cm}p{7.8cm}@{}}
\caption{The 49 time-domain descriptors, in the order they appear in the
feature bank for each channel.}
\label{tab:tddefs}\\
\toprule
\textbf{\#} & \textbf{Name} & \textbf{Definition} \\
\midrule
\endfirsthead
\toprule
\textbf{\#} & \textbf{Name} & \textbf{Definition} \\
\midrule
\endhead
1 & mean & Arithmetic mean. \\
2 & median & Median. \\
3 & mode & Most frequent value. \\
4 & std & Standard deviation. \\
5 & variance & Variance. \\
6 & skewness & Sample skewness. \\
7 & kurtosis & Pearson (non-excess) kurtosis. \\
8 & coeff.\ of variation & std / mean (tsfresh \texttt{variation\_coefficient}). \\
9 & frac.\ zero values & Fraction of samples exactly equal to zero (tsfresh \texttt{value\_count}). \\
10 & frac.\ NaNs & Fraction of NaN samples in the raw window. \\
11 & longest strike above mean & Longest consecutive run above the mean. \\
12 & longest strike below mean & Longest consecutive run below the mean. \\
13 & number of peaks (lag) & Count of local maxima with the given lag as minimum peak support. \\
14 & sum of reoccurring values & Sum of values that occur more than once in the window. \\
15 & count of zeros & Same tsfresh calculator as \#9; a duplicate column. \\
16 & pct.\ reoccurring points & Percentage of samples that are non-unique. \\
17 & ratio unique values & Unique-value count / window length. \\
18 & autocorrelation (lag) & Autocorrelation at the window's lag. \\
19 & partial autocorr.\ (lag) & Partial autocorrelation at the window's lag. \\
20 & Benford correlation & Correlation between the leading-digit distribution and Benford's law. \\
21 & c3 nonlinearity & tsfresh \texttt{c3} nonlinearity statistic at the window's lag. \\
22 & cid-ce complexity & Disabled in the released code; identically zero. \\
23 & symmetry-looking flag & Binary: distribution symmetric about the mean within $r=0.01$. \\
24 & time-reversal asymmetry & tsfresh \texttt{time\_reversal\_asymmetry\_statistic} at the window's lag. \\
25 & RMS & Root-mean-square. \\
26 & absolute energy & Disabled in the released code; identically zero. \\
27 & absolute sum of changes & Disabled in the released code; identically zero. \\
28 & mean abs.\ change & Mean of $|x_{t+1}-x_t|$. \\
29 & mean change & Mean of $x_{t+1}-x_t$. \\
30 & mean 2nd derivative & Mean of a central second-difference estimate. \\
31 & variance$>$std flag & Binary: variance exceeds $1$. \\
32 & large-std flag & Binary: std exceeds $25\%$ of the range. \\
33 & ratio beyond 1-sigma & Fraction of samples more than $1\sigma$ from the mean. \\
34 & var.\ of STA histogram & std/mean of a $26$-bin histogram of the $1$-s moving-average trace. \\
35 & var.\ of MTA histogram & Same, $20$-s moving average. \\
36 & var.\ of LTA histogram & Same, $120$-s moving average. \\
37 & homogeneity of STA hist.\ & max/mean of the $1$-s moving-average histogram. \\
38 & homogeneity of MTA hist.\ & Same, $20$-s moving average. \\
39 & homogeneity of LTA hist.\ & Same, $120$-s moving average. \\
40 & PSD entropy & Disabled in the released code; identically zero. \\
41 & FFT entropy & Disabled in the released code; identically zero. \\
42 & PSD kurtosis & Kurtosis of the normalized windowed-PSD estimate of the channel itself. \\
43 & FFT kurtosis & Kurtosis of the normalized whole-window FFT magnitude-squared. \\
44 & PSD homogeneity & max/mean of the normalized windowed-PSD estimate. \\
45 & FFT homogeneity & max/mean of the normalized whole-window FFT magnitude-squared. \\
46 & PSD bandwidth (95pct) & Frequency range containing $95\%$ of the windowed-PSD energy. \\
47 & PSD bandwidth (50pct) & Frequency range containing $50\%$ of the windowed-PSD energy. \\
48 & FFT bandwidth (95pct) & Same, whole-window FFT, $95\%$. \\
49 & FFT bandwidth (50pct) & Same, whole-window FFT, $50\%$. \\
\bottomrule
\end{longtable}

\subsection{PSD and transmissibility descriptors (74 per channel or pair)}

The same $74$ descriptors are computed from two kinds of $300$-point spectra
spanning $0$--$30$~Hz at $0.1$~Hz resolution: the per-channel PSD (Welch's
method, \texttt{nperseg}$=1000$ at $100$~Hz) and the transmissibility ratio
$T_{ij}(f)$ of Eq.~4 of the main text. The $74$ values fall into five
groups:

\begin{itemize}
\item \textbf{Peak and amplitude statistics (1--10):} mean amplitude; number
of peaks (tsfresh \texttt{number\_peaks}, minimum separation $25$ bins);
standard deviation, mean, kurtosis, and skewness of the peak positions;
mean and standard deviation of peak prominence; mean and standard deviation
of peak height.
\item \textbf{Piecewise-cubic regression coefficients (11--34, $24$ values):}
the spectrum is split into six non-overlapping $50$-bin ($5$~Hz) segments;
a degree-$3$ polynomial is fit to each, retaining all four coefficients
(cubic, quadratic, linear, constant), giving $6\times4=24$ values.
\item \textbf{Sharpest-peak indices, positions, and heights (35--46, $12$
values):} among all detected peaks, the four with the largest prominence
are ranked; their index rank, spectral position, and height are recorded
($4$ values each).
\item \textbf{Sub-band peak height and relative position (47--70, $24$
values):} the spectrum is split into twelve $2.5$~Hz sub-bands; the peak
value in each sub-band and its position within the sub-band (as a
percentage) are recorded ($12$ values each).
\item \textbf{Spectral centroid (71--74, $4$ values):} the spectrum is split
into four $7.5$~Hz quarters; the frequency-weighted centroid of each
quarter is recorded.
\end{itemize}

\begin{thebibliography}{1}
\bibitem{christ2018}
M.~Christ, N.~Braun, J.~Neuffer, A.~W. Kempa-Liehr,
``Time series feature extraction on basis of scalable hypothesis tests
(tsfresh -- a Python package),''
\emph{Neurocomputing}, 307:72--77, 2018.
\url{https://doi.org/10.1016/j.neucom.2018.03.067}
\end{thebibliography}

\end{document}
